\documentclass[10pt]{article}
\usepackage[margin=0.55in]{geometry}
\usepackage{amsmath, amssymb}
\usepackage{enumitem}
\usepackage{titlesec}
\usepackage{array}
\setlist{nosep, leftmargin=1.2em}
\titlespacing*{\section}{0pt}{5pt}{2pt}
\titleformat{\section}{\normalsize\bfseries}{}{0pt}{}
\setlength{\parskip}{3pt}
\setlength{\parindent}{0pt}
\pagestyle{empty}

\begin{document}

\begin{center}{\large\bfseries ACI 318-19 Punching-Shear Design Conventions}\end{center}

One toggle in the Inputs panel turns on two ACI-consistent design conventions at once. Both are long-standing ACI provisions; both reduce false positives produced by a raw biaxial check on an unstable solver or a symmetric tributary. The flow is \textit{first} try to qualify the column for $\gamma_v = 0$ via Table 8.4.2.2.4; \textit{then}, at corner columns only, keep $0.3\,M_o$ as a floor so a solver can't silently under-predict at the column type most sensitive to live-load patterning.

\section*{1. Modified $\gamma_f$ per Table 8.4.2.2.4}
At a column the unbalanced slab moment $M_{sc}$ is split: a fraction $\gamma_f M_{sc}$ is transferred by flexure across the effective strip $c_2 + 3h$, and $\gamma_v M_{sc} = (1-\gamma_f) M_{sc}$ by eccentric shear on the critical section. Baseline (\S 8.4.2.2.2) gives $\gamma_f = 1/\bigl(1 + \tfrac{2}{3}\sqrt{b_1/b_2}\bigr)$. Table 8.4.2.2.4 permits $\gamma_f$ to be raised --- up to $1.0$ for corner and edge-perpendicular cases --- when the direct shear is modest AND the slab rebar is ductile.

\begin{tabular}{|l|l|c|c|c|}
\hline
\textbf{Column} & \textbf{Span direction} & $v_{uv} \leq$ & $\varepsilon_t \geq$ & $\gamma_{f,\max}$ \\
\hline
Corner          & either                   & $0.5\,\phi v_c$   & $\varepsilon_{ty} + 0.003$ & $1.0$ \\
Edge            & perpendicular to edge    & $0.75\,\phi v_c$  & $\varepsilon_{ty} + 0.003$ & $1.0$ \\
Edge            & parallel to edge         & $0.4\,\phi v_c$   & $\varepsilon_{ty} + 0.008$ & $\dfrac{1.25}{1 + \tfrac{2}{3}\sqrt{b_1/b_2}} \leq 1.0$ \\
Interior        & either                   & $0.4\,\phi v_c$   & $\varepsilon_{ty} + 0.008$ & $\dfrac{1.25}{1 + \tfrac{2}{3}\sqrt{b_1/b_2}} \leq 1.0$ \\
\hline
\end{tabular}

\vspace{3pt}

\noindent\textbf{(a) Direct shear.} $v_{uv} = V_u / (b_0 d)$, excluding moment contributions. Must fall at or below the tabulated fraction of $\phi v_c$ for the row to apply. The app computes $v_{uv}$ per column and falls back to baseline $\gamma_f$ when the gate fails --- no user action required.

\vspace{2pt}

\noindent\textbf{(b) Ductile flexure.} $\varepsilon_t$ in the column-strip reinforcement must meet the stated strain minimum. For Grade 60 these evaluate to $\approx 0.00507$ (rows with $\gamma_{f,\max} = 1.0$) and $\approx 0.01007$ (rows with the $1.25\times$ boost). Typical well-designed slabs satisfy the looser bound without trying; the stricter bound pushes toward lightly-reinforced column strips. \textbf{Checking this box is your attestation that the slab design meets this limit and has rebar concentrated across $c_2 + 3h$ per \S 8.4.2.2.5}.

\section*{2. $0.3\,M_o$ Floor per DDM (\S 8.10) --- corner columns only}
After $\gamma_f$ is fixed, a lower bound is enforced on $|M_u|$ \textbf{at corner columns only}:
\[
|M_{u,\text{used}}|
\;=\; \max\!\bigl(\,|M_{u,\text{solver}}|,\, 0.3\, M_o\,\bigr), \qquad
M_o \;=\; \frac{w_u\, l_2\, l_n^2}{8}.
\]
$l_1$ and $l_2$ are center-to-center spans to adjacent columns (or to the slab edge at a free side); $l_n \geq 0.65\,l_1$. The floor is applied per axis independently. Sign is preserved from the solver. Interior and edge columns use the solver's $M_u$ unchanged.

\vspace{3pt}

\noindent\textbf{Why corner-only.} Corners have the smallest tributary and the largest relative sensitivity to live-load patterning --- a one-sided pattern that would nudge an interior Mu by a few percent can dominate a corner's Mu, and gravity-balanced FEA will not see that. DDM's $0.3\,M_o$ is a historically accepted lower bound for exactly this pattern-loading case, and keeping the floor narrow to corner columns avoids over-penalizing edge/interior columns where a stable FEA already captures asymmetric tributary correctly. Using it as a floor (not a ceiling) keeps solver detail where it exceeds DDM and patches solver drift where it doesn't.

\vspace{2pt}

\noindent\textbf{Interaction.} When the Table 8.4.2.2.4 gate passes for an axis, $\gamma_v = 0$ on that axis, and the eccentric-shear term drops out regardless of $M_u$. In that case the $0.3\,M_o$ floor does not affect the DCR for that axis --- but the flexural design at the column strip still has to carry $\gamma_f \cdot M_u$, and that's still driven by $\max(\text{solver}, 0.3 M_o)$.

\section*{What the toggle does NOT do}
\begin{itemize}
\item Does not verify the rebar detailing. Checking the box is an attestation on your part.
\item Does not replace the solver. Solver Mu always wins when it exceeds the floor.
\item Does not redistribute between axes. The floor is applied per column per axis.
\item Does not apply to seismic or lateral-driven Mu --- those have their own code provisions and should be input directly as overrides.
\end{itemize}

\section*{When to leave it off}
Turn it off (i) when reviewing results against a raw biaxial ACI \S 8.4.2.2.2 check with no DDM allowances, (ii) when the slab design is not tension-controlled or the column-strip concentration per \S 8.4.2.2.5 is not provided, or (iii) when the governing condition is non-gravity and $0.3\,M_o$ does not represent the lower bound that applies.

\end{document}
